% GENERATED FILE. Edit canonical lessons, then run npm run generate:books.
% canonical-source-hash: fnv1a64:17972da21091a936
% canonical-lessons: MR-C137-pohe, MR-C137-ladu, MR-C137-kakdi, MR-C137-lasun, MR-C137-ale

\chapter{Flattened rice, Ladoo, Cucumber, Garlic, Ginger}
\label{ch:mr-flattened-rice-ladoo-cucumber-garlic-ginger}
\hlchaptermodality{137}

\begin{chapteropening}
\textbf{By the end of this chapter:} \emph{I can say flattened rice, a ladoo, a cucumber, garlic and ginger.}

Complete the last lesson of chapter 137: I can say flattened rice, a ladoo, a cucumber, garlic and ginger.
\end{chapteropening}

\section[\texorpdfstring{\mr{पोहे} (pohe)}{पोहे (pohe)}]{\mr{पोहे} (pohe) — flattened rice}

\label{lesson:MR-C137-pohe}



\begin{quote}
Before the new one: say the Marathi for a vegetable dish, then the Marathi for plain dal.
\end{quote}

\subsection*{You'll want to know: \mr{पोहे}}
\textbf{\mr{पोहे}} — \emph{pohe} — \textquotedblleft{}flattened rice\textquotedblright{}.

\begin{cousinweb}
\textbf{\mr{पोहे}} (\emph{pohe}) is a tadbhava, a word inherited from Sanskrit and reshaped by sound change. It comes from \emph{pṛthuka}, \textquotedblleft{}flattened rice or grain\textquotedblright{} (CDIAL 8359). It is built on \emph{pṛthu}, \textquotedblleft{}broad, flat\textquotedblright{}.
\end{cousinweb}

\begin{grammarlens}[title={What you've built}]
One more word for this chapter.
\end{grammarlens}

\subsection*{Your turn}
\begin{itemize}
\raggedright
  \item \emph{Say it:} \emph{pohe}
  \item \emph{Say it:} \emph{pohe}, once more
  \item \emph{Say it:} say \emph{pohe}
  \item \emph{Recall:} say the Marathi for a vegetable dish, then the Marathi for plain dal, then say \emph{pohe} again
\end{itemize}

\begin{tcolorbox}[breakable,colback=teal!4,colframe=teal!35!black,title={Before you move on}]
What does \mr{पोहे} mean? (\textquotedblleft{}Flattened rice\textquotedblright{}.) Say it once more.
\end{tcolorbox}
\section[\texorpdfstring{\mr{लाडू} (lāḍū)}{लाडू (lāḍū)}]{\mr{लाडू} (lāḍū) — a ladoo, a sweet ball}

\label{lesson:MR-C137-ladu}



\begin{quote}
Before the new one: say the Marathi for plain dal, then the Marathi for flattened rice.
\end{quote}

\subsection*{You'll want to know: \mr{लाडू}}
\textbf{\mr{लाडू}} — \emph{lāḍū} — \textquotedblleft{}a ladoo, a sweet ball\textquotedblright{}.

\begin{cousinweb}
\textbf{\mr{लाडू}} (\emph{lāḍū}) is a tadbhava of Sanskrit \emph{laḍḍu}, \textquotedblleft{}a kind of sweetmeat\textquotedblright{} (CDIAL 10926). CDIAL takes it as probably the same word as \emph{*laḍḍu}, \textquotedblleft{}lump\textquotedblright{}, and notes Kuiper's link with Munda words for \textquotedblleft{}lump\textquotedblright{}.
\end{cousinweb}

\begin{grammarlens}[title={What you've built}]
One more word for this chapter.
\end{grammarlens}

\subsection*{Your turn}
\begin{itemize}
\raggedright
  \item \emph{Say it:} \emph{lāḍū}
  \item \emph{Say it:} \emph{lāḍū}, once more
  \item \emph{Say it:} say \emph{lāḍū}
  \item \emph{Recall:} say the Marathi for plain dal, then the Marathi for flattened rice, then say \emph{lāḍū} again
\end{itemize}

\begin{tcolorbox}[breakable,colback=teal!4,colframe=teal!35!black,title={Before you move on}]
What does \mr{लाडू} mean? (\textquotedblleft{}A ladoo, a sweet ball\textquotedblright{}.) Say it once more.
\end{tcolorbox}
\section[\texorpdfstring{\mr{काकडी} (kākḍī)}{काकडी (kākḍī)}]{\mr{काकडी} (kākḍī) — a cucumber}

\label{lesson:MR-C137-kakdi}



\begin{quote}
Before the new one: say the Marathi for flattened rice, then the Marathi for a ladoo.
\end{quote}

\subsection*{You'll want to know: \mr{काकडी}}
\textbf{\mr{काकडी}} — \emph{kākḍī} — \textquotedblleft{}a cucumber\textquotedblright{}.

\begin{cousinweb}
\textbf{\mr{काकडी}} (\emph{kākḍī}) is a tadbhava, a word inherited from Sanskrit and reshaped by sound change. It comes from \emph{karkaṭī}, \textquotedblleft{}Cucumis utilissimus, a cucumber\textquotedblright{} (CDIAL 2817).
\end{cousinweb}

\begin{grammarlens}[title={What you've built}]
One more word for this chapter.
\end{grammarlens}

\subsection*{Your turn}
\begin{itemize}
\raggedright
  \item \emph{Say it:} \emph{kākḍī}
  \item \emph{Say it:} \emph{kākḍī}, once more
  \item \emph{Say it:} say \emph{kākḍī}
  \item \emph{Recall:} say the Marathi for flattened rice, then the Marathi for a ladoo, then say \emph{kākḍī} again
\end{itemize}

\begin{tcolorbox}[breakable,colback=teal!4,colframe=teal!35!black,title={Before you move on}]
What does \mr{काकडी} mean? (\textquotedblleft{}A cucumber\textquotedblright{}.) Say it once more.
\end{tcolorbox}
\section[\texorpdfstring{\mr{लसूण} (lasūṇ)}{लसूण (lasūṇ)}]{\mr{लसूण} (lasūṇ) — garlic}

\label{lesson:MR-C137-lasun}



\begin{quote}
Before the new one: say the Marathi for a ladoo, then the Marathi for a cucumber.
\end{quote}

\subsection*{You'll want to know: \mr{लसूण}}
\textbf{\mr{लसूण}} — \emph{lasūṇ} — \textquotedblleft{}garlic\textquotedblright{}.

\begin{cousinweb}
\textbf{\mr{लसूण}} (\emph{lasūṇ}) is a tadbhava, a word inherited from Sanskrit and reshaped by sound change. It comes from \emph{laśuna}, \textquotedblleft{}garlic\textquotedblright{} (CDIAL 10990).
\end{cousinweb}

\begin{grammarlens}[title={What you've built}]
One more word for this chapter.
\end{grammarlens}

\subsection*{Your turn}
\begin{itemize}
\raggedright
  \item \emph{Say it:} \emph{lasūṇ}
  \item \emph{Say it:} \emph{lasūṇ}, once more
  \item \emph{Say it:} say \emph{lasūṇ}
  \item \emph{Recall:} say the Marathi for a ladoo, then the Marathi for a cucumber, then say \emph{lasūṇ} again
\end{itemize}

\begin{tcolorbox}[breakable,colback=teal!4,colframe=teal!35!black,title={Before you move on}]
What does \mr{लसूण} mean? (\textquotedblleft{}Garlic\textquotedblright{}.) Say it once more.
\end{tcolorbox}
\section[\texorpdfstring{\mr{आले} (āle)}{आले (āle)}]{\mr{आले} (āle) — ginger}

\label{lesson:MR-C137-ale}



\begin{quote}
Before the new one: say the Marathi for a cucumber, then the Marathi for garlic.
\end{quote}

\subsection*{You'll want to know: \mr{आले}}
\textbf{\mr{आले}} — \emph{āle} — \textquotedblleft{}ginger\textquotedblright{}.

\begin{cousinweb}
\textbf{\mr{आले}} (\emph{āle}) is a tadbhava, a word inherited and reshaped by sound change. It comes from \emph{*āllaka}, \textquotedblleft{}fresh ginger\textquotedblright{}, a form CDIAL reconstructs from its descendants (CDIAL 1341). CDIAL builds it from Sanskrit \emph{ārdraka}, \textquotedblleft{}fresh ginger\textquotedblright{}.
\end{cousinweb}

\begin{grammarlens}[title={What you've built}]
One more word for this chapter.
\end{grammarlens}

\subsection*{Your turn}
\begin{itemize}
\raggedright
  \item \emph{Say it:} \emph{āle}
  \item \emph{Say it:} \emph{āle}, once more
  \item \emph{Say it:} say \emph{āle}
  \item \emph{Recall:} say the Marathi for a cucumber, then the Marathi for garlic, then say \emph{āle} again
\end{itemize}

\begin{tcolorbox}[breakable,colback=teal!4,colframe=teal!35!black,title={Before you move on}]
What does \mr{आले} mean? (\textquotedblleft{}Ginger\textquotedblright{}.) Say it once more.
\end{tcolorbox}
